\par\clearpage
\noindent\null\hfill{\small\textbf{ZUR HANNOVERSCHEN TRIANGULATION.}}\hfill\makebox[0pt][r]{\fontsize{13}{16}\selectfont 411}\par\medskip
\noindent sich ferner die Möglichkeit, den Falkenberg mit Wilsede mittels eines
Durchhaus durch das Becklinger Holz zu verbinden, und die Hoffnung, dass
die beiden Punkte Hauselberg und Garssen mittels eines oder einiger
Zwischenpunkte und einiger Durchhaue durch die Waldungen, die sie trennen,
verknüpft werden könnten. Hinsichtlich solcher Durchhaue jedoch, ohne die
in diesen Gegenden schlechterdings kein Durchkommen ist, glaubte ich mir
zwei Regeln auferlegen zu müssen: erstens, soweit irgend möglich, allen
edleren Holzarten (etwa Eichenwäldern) auszuweichen, und zweitens die Richtung
der Durchhaue \textls[120]{vorher} mit \textls[120]{äußerster Präzision} zu
bestimmen, damit sie so schmal wie möglich ausfielen und kein einziger Stamm
ohne Not gefällt werden musste. Um das zu erreichen, mussten jedoch bereits
sehr genaue Messungen vorausgehen, die zum Teil geschickt angeordnet und
kombiniert werden mussten. Diese ließen sich erst dann mit Genauigkeit und
mit dem geringstmöglichen Zeitaufwand ausführen, wenn die Hauptmessungen
selbst bis in diese Gegend vorgerückt waren. Ich darf hier vorweg bemerken,
dass ich diese Ziele später auch zu meiner vollkommensten Zufriedenheit
erreicht habe.

Die Zeit zwischen meiner Rückkehr von der Erkundungsreise und dem eigentlichen
Beginn der Messungsarbeiten nutzte ich, um die Instrumente ganz gebrauchsfertig
zu machen, alles Nötige vorzubereiten und dem Hauptmann \textsc{Müller} und
dem Leutnant \textsc{Hartmann} die Anweisungen zu geben, die für ein vollkommenes
Ineinandergreifen der Arbeiten erforderlich waren. Am 16.~Juni trat ich die
Reise nach Lichtenberg an, dem ersten in diesem Jahr zu bearbeitenden
Dreieckspunkt, wohin mein Sohn mit einem Teil der Instrumente bereits
vorausgereist war.

Die eigentlichen Messungsarbeiten haben in diesem Jahr vier Monate gedauert,
nämlich bis Mitte Oktober. Während dieser Zeit wurden die Arbeiten an
\textls[120]{neun} Hauptdreieckspunkten abgeschlossen: Im Jahr 1821 konnten in einer
nicht viel kürzeren Zeit nur fünf bearbeitet werden. Dieser so viel raschere
Fortschritt ist hauptsächlich der größeren Zahl der Gehilfen und der
Heliotrope zuzuschreiben. Noch wichtiger ist aber die dadurch erreichte
Steigerung der Genauigkeit der Messungen selbst.\quad\mbox{.\ .\ .\ .\ .\ .}

Der Zweck der Triangulation als Teil der Gradmessung ist, die Göttinger
Sternwarte durch ein zusammenhängendes System von Dreiecken mit den \mbox{däni-}
\par\smallskip\noindent\hfill\makebox[0.18\linewidth][l]{52\textsuperscript{*}}\par
\clearpage
\noindent\makebox[0pt][l]{\fontsize{13}{16}\selectfont 412}\hfill{\small\textbf{NACHLASS.}}\hfill\null\par\medskip
\noindent schen Dreiecken zu verbinden. Dazu war es am vorteilhaftesten, die Dreiecke
so groß wie möglich anzulegen und die Dreieckspunkte im Allgemeinen an den
höchsten Stellen zu wählen, die die weiteste Aussicht gewähren. Diese Punkte
sind an und für sich größtenteils ohne unmittelbares Interesse für die
Geographie des Königreichs. Neben dem Hauptzweck habe ich aber überall
versucht, meine Arbeiten für diese Geographie möglichst nützlich zu machen.
Die Lage aller Orte, die von mehr als einem Hauptdreieckspunkt aus sichtbar
waren, habe ich sorgfältig bestimmt, manche mit der gleichen Genauigkeit wie
die Lage der Hauptdreieckspunkte. Ich nenne davon die Städte Hannover,
Braunschweig, Celle, Lüneburg, Neustadt am Rübenberge und Burgdorf. Die Zahl
der Dörfer, deren Lage ich bestimmt habe, ist sehr groß. Der Weg ist gebahnt,
diese Ernte auf einen größeren Teil des Königreichs oder auf das ganze
Königreich auszudehnen.\quad\mbox{.\ .\ .\ .\ .\ .}

\bigskip
\begin{center}
[Aus dem Bericht vom 16.~Februar 1825 ȟber die im Jahre 1824\par
ausgeführten trigonometrischen Arbeiten.«]
\end{center}

\par\noindent\hspace*{\parindent}\mbox{.\ .\ .\ .\ .\ .}\quad Resultate: Der Bericht und die beigefügte Karte zeigen, dass
die Arbeiten des Jahres 1824 mehrere Verbindungen von den Dreiecken der
früheren Jahre bis Bremen bieten. Die einfachste [die Dreiecke
Wilsede-Falkenberg-Elmhorst, Wilsede-Elmhorst-Litberg, Wilsede-Litberg-Hamburg,
Wilsede-Litberg-Zeven, Wilsede-Zeven-Steinberg und Zeven-Steinberg-Bremen
umfassend\,*)] ist mit starken durchgezogenen Linien gezeichnet, und zwei neue
Dreiecke [Zeven-Bremen-Brillit und Brillit-Bremen-Garlste] sind noch an die
Seite Bremen-Zeven angeknüpft. Wäre es möglich gewesen, jenes einfachste
System gleich anfangs ausfindig zu machen, so hätten allerdings die anderen
mit dünnen durchgezogenen Linien gezeichneten Dreiecke [bei denen die Punkte
Bullerberg, Bottel und Brüttendorf Eckpunkte sind] ganz wegfallen können. Der
Bericht zeigt jedoch, wie viele Schwierigkeiten zu überwinden waren, ehe der
Plan dafür ermittelt werden konnte, und die präzise und schnelle Ausführung
der verschiedenen dazu erforderlichen Durchhaue wäre ebenfalls ohne vorherige
\par\vfill\noindent\rule{0.25\linewidth}{0.4pt}\par\smallskip
{\footnotesize\noindent\hspace*{\parindent}[*) Siehe die Dreiecksskizze auf S.~299].\par}
\clearpage
\noindent\null\hfill{\small\textbf{ZUR HANNOVERSCHEN TRIANGULATION.}}\hfill\makebox[0pt][r]{\fontsize{13}{16}\selectfont 413}\par\medskip
\noindent sehr genaue Kenntnis der Lage der Plätze ganz undurchführbar gewesen.
Beim heutigen mathematischen Stand der höheren Geodäsie dürfen übrigens auch
die letzteren Dreiecke, die dünn gezeichneten, keineswegs als überflüssig
betrachtet werden. Vielmehr muss ihre Berücksichtigung nach ganz bestimmten
Grundsätzen dazu beitragen, die Genauigkeit der Endresultate zu erhöhen.\quad\mbox{.\ .\ .\ .\ .\ .}

Schließlich ist auch noch von großer Bedeutung, dass durch die drei Dreiecke
zwischen den fünf Punkten Falkenberg, Elmhorst, Wilsede, Litberg, Hamburg
eine neue Verbindung von den südlichen Dreiecken im Königreich Hannover bis
Hamburg erreicht worden ist, die der früheren, viel komplizierteren [über
Hauselberg, Wulfsode, Timpenberg, Nindorf und Lüneburg] vorzuziehen ist und
daher nach den zuvor angedeuteten Grundsätzen die Genauigkeit der Resultate
verdoppeln wird.\quad\mbox{.\ .\ .\ .\ .\ .}

\bigskip
\begin{center}
[Aus dem Bericht vom 21.~November 1827, »betreffend die weitere Ausdehnung\par
der Gradmessungsarbeiten.«]
\end{center}

\par\noindent\hspace*{\parindent}\mbox{.\ .\ .\ .\ .\ .}\quad Für die Messung der Dreieckskette bis Hamburg wurden
1821---1823 11000 Taler ausgegeben und 1824, 1825 für die von dort westlich bis
Ostfriesland geführte Dreieckskette etwa 7000 Taler. Von diesen Kosten ist
aber abzuziehen, was für die Anschaffung von Instrumenten und für das Abholen
des englischen Zenitsektors von Altona nach Göttingen ausgegeben worden ist
und zwischen 2500 und 3000 Talern betragen haben mag, sodass die eigentlichen
Triangulationskosten etwa 15000 Taler betragen haben mögen. Nach der
Übersichtskarte scheint nun die Gesamtheit der noch nicht erfassten Landesteile
wohl nicht viel größer zu sein als die bereits mit Dreiecken überzogene Fläche.
Bei aller Ungewissheit über die Schwierigkeiten des Geländes ist es doch kaum
wahrscheinlich, dass sie größer sein können als diejenigen, mit denen ich
besonders 1822 und 1824 zu kämpfen hatte. Wenn ich nun außerdem bemerke, dass
die Arbeiten, deren Hauptzweck die Vervollkommnung der Landesgeographie ist,
auch bei einer angemessenen Ausführung doch nicht \textls[120]{den} Grad
äußerster Messgenauigkeit erfordern, der bei einer eigentlichen Gradmessung
verlangt wird, so scheint
\par\clearpage
\noindent\makebox[0pt][l]{\fontsize{13}{16}\selectfont 414}\hfill{\small\textbf{NACHLASS.}}\hfill\null\par\medskip
\noindent die Hoffnung nicht unbegründet, dass die Erweiterung der Triangulation
auf die noch nicht erfassten Teile des Königreichs mit einer geringeren
Summe und vielleicht mit 12000 Talern zu bestreiten sein wird\,*).

Um eine solche Triangulation für die Vervollkommnung der Geographie möglichst
nützlich zu machen, wird man sich nicht darauf beschränken dürfen, lediglich
Netze von Hauptdreiecken erster Ordnung über die betreffenden Landesteile
zu legen, sondern damit die Lagebestimmung einer möglichst großen Zahl
sekundärer Punkte verbinden müssen, namentlich solcher, die genaue
Bestimmungen zulassen und in der Regel Jahrhunderte überdauern, also
besonders der Kirchtürme. Diese Rücksicht habe ich mir schon bei den früheren
Messungen zur Pflicht gemacht, obwohl sie dem Hauptzweck untergeordnet
bleiben musste. Die Zahl der bei jenen Messungen bestimmten Punkte beträgt
schon über 500.\quad\mbox{.\ .\ .\ .\ .\ .}

Diese Angabe der Lage einer großen Zahl fester Punkte in Zahlen (also wie
weit nördlich oder südlich, westlich oder östlich von einem beliebigen
Ausgangspunkt, z.~B. der Göttinger Sternwarte), \textls[120]{bis auf wenige
Fuß genau}, muss in topographischer Hinsicht als Hauptertrag der Arbeiten
betrachtet werden. Sie behält auf Jahrhunderte ihren Wert, da die Mehrzahl der
Punkte bestehen bleibt, auch wenn im Laufe der Zeit einige verloren gehen.
Die dadurch etwa entstehenden Veränderungen lassen sich leicht ergänzen.
Sie bildet eine sichere Grundlage für alle Detailaufnahmen: Alle
Unsicherheiten, die Aufnahmen ohne solche festen Anhaltspunkte erschweren,
entstellen und ihre Vereinigung zu einem fehlerfreien Ganzen unmöglich
machen, fallen dabei ganz weg. Nachlässige Arbeiter unterliegen dadurch
einer strengen und unausweichlichen Kontrolle. Jede Messtischplatte wird
unabhängig von der anderen bearbeitet, sodass sich kein Fehler auf andere
Blätter fortpflanzt. Schließlich fügen sich alle einzelnen Blätter von selbst
zu einem genau orientierten und überall zusammenpassenden Ganzen. Es ist
einleuchtend, dass die hohen Kosten, die Detailaufnahmen von bedeutendem
Umfang immer verursachen, durch ein solches sicheres Vorgehen in hohem Maße
vermindert werden müssen. Doch ist es nicht allein dieses sichere Vorgehen,
\par\vfill\noindent\rule{0.25\linewidth}{0.4pt}\par\smallskip
{\footnotesize\noindent\hspace*{\parindent}*) Die im laufenden Jahr 1827 erforderlichen Kosten bleiben hierbei außer Betracht, da ihr Gegenstand zum rein astronomischen Teil der Gradmessung gehört.\par}
\clearpage
\noindent\null\hfill{\small\textbf{ZUR HANNOVERSCHEN TRIANGULATION.}}\hfill\makebox[0pt][r]{\fontsize{13}{16}\selectfont 415}\par\medskip
\noindent das die Arbeit beschleunigt. In dieser Hinsicht ist es auch sehr wichtig,
dass der Gebrauch der Messkette dadurch fast ganz überflüssig und nur
ausnahmsweise nötig wird, da die Triangulation die Grundlinien bereits von
selbst liefert, und zwar mit einer Genauigkeit, die die gewöhnliche Kette
gar nicht erreichen könnte.

Aber auch dort, wo bereits Detailaufnahmen vorliegen, wie bei den meisten
Ämtern des früheren Bestands des Königreichs, bieten die festen Punkte ein
Mittel, die beim Zusammensetzen entstandenen Fehler zu berichtigen. So lassen
sich selbst Karten, die auf unvollkommenen Aufnahmemethoden beruhen, wenn
sie sonst im Detail gut sind, zu Darstellungen umarbeiten, die auch höheren
Anforderungen genügen können.

Was demnach die Maßnahmen betrifft, um die Triangulationen für die
Vervollkommnung der Geographie möglichst nützlich zu machen, so sind die
bereits ausgeführten Messungen von denen zu unterscheiden, die künftig
eventuell auf andere Landesteile ausgedehnt werden.

Bei letzteren wird die genaue Lagebestimmung einer möglichst großen Zahl
fester Punkte gleich als Hauptzweck berücksichtigt werden müssen.

Bei den bereits ausgeführten Messungen hingegen ist diese Rücksicht
allerdings nur als untergeordnet betrachtet worden. Da ich sie jedoch, wie
schon erwähnt, stets im Auge behalten habe, soweit dies ohne Beeinträchtigung
der Hauptaufgabe möglich war, müssen hinsichtlich des Ergebnisses hier
erneut die nördlichen von den südlichen Gegenden unterschieden werden.

In der nördlichen (größeren) Hälfte, d.~h. etwa von der Stadt Hildesheim bis
zum Meer, also im flachen Teil des Landes, ist der Ertrag in der genannten
Hinsicht so reich gewesen, dass wenig oder nichts zu wünschen übrig bleibt.\quad\mbox{.\ .\ .\ .\ .\ .}

Im südlichsten Teil des Königreichs hingegen ist die Zahl der genau
bestimmten Kirchtürme viel kleiner, da teils wegen der Größe der Dreiecke,
teils wegen der gebirgigen Beschaffenheit des Landes nur wenige Türme von
mehr als einem Hauptdreieckspunkt aus zugleich sichtbar waren. Für die
Vervollkommnung der Geographie des Königreichs und insbesondere, um einer
Detailaufnahme der südlichen Teile des Hildesheimschen und des Eichsfeldes
ähnlich sichere Grundlagen zu verschaffen, wäre es daher \mbox{aller-}
\par
